%% kapa: exam=1 round=27 year=2016 period=1 form=A subject=economics subject_ko=경제학원론 seq=5 number=45
%% source: Q-Net 감정평가사 기출문제 (cst00309) artlSeq=5169402 file=extracted/1008/5169402/제27회 감정평가사 1차 문제/제27회 감정평가사 1교시문제(A형).pdf
%% conversion: kordoc 4.24.51 -> tools/gen1.py (auto); review=proofread-sonnet; flags=none
\begin{question}{45}
甲의 효용함수는 U(x, y) = xy이고, X재와 Y재의 가격이 각각 2,000원과 8,000원이며, 소득은 100,000원이다. 예산제약 하에서 甲의 효용이 극대화되는 소비점에서 한계대체율($MRS_{XY} = -\Delta Y/\Delta X$)은? (단, 甲은 X재와 Y재만 소비하고, x는 X재의 소비량, y는 Y재의 소비량이다.)
\begin{choices}
  \item 0.25
  \item 0.5
  \item 0.75
  \item 2.0
  \item 2.5
\end{choices}
\end{question}
\answer{1}
